% Électrolyse de l'eau : éprouvettes graduées (exercice 6, chapitre 10)
\begin{tikzpicture}[x=1cm,y=1cm]
\pic at (0,0) {becher={largeur=5.2,hauteur=2.6,niveau=2.15,liquide=pcEau}};
\foreach \x in {-1.25,1.25}{%
  \fill[pcEau] (\x-0.3,0.6) rectangle (\x+0.3,3.8);}
\solution{\fill[white] (-1.55,2.96) rectangle (-0.95,3.8);
  \draw[pcSurface] (-1.55,2.96) -- (-0.95,2.96);
  \fill[white] (0.95,3.38) rectangle (1.55,3.8);
  \draw[pcSurface] (0.95,3.38) -- (1.55,3.38);
  \node[pcLegende,anchor=east,text=otGreenDark] at (-1.75,3.35) {28~mL};
  \node[pcLegende,anchor=west,text=otGreenDark] at (1.75,3.55) {14~mL};}
\foreach \x in {-1.25,1.25}{%
  \draw[pcVerre] (\x-0.36,0.55) -- (\x-0.36,3.74) arc[start angle=180,end angle=0,radius=0.36] -- (\x+0.36,0.55);
  \foreach \k in {0,...,6}{\draw[otNavy!50,line width=0.4pt] (\x+0.21,{1.5+\k*0.37}) -- (\x+0.35,{1.5+\k*0.37});}}
\filldraw[pcAppareil] (-1.37,0.32) rectangle (-1.13,1.3);
\filldraw[pcAppareil] (1.13,0.32) rectangle (1.37,1.3);
\draw[pcFil] (-1.25,0.32) -- (-1.25,-0.45) -- (-0.45,-0.45);
\draw[pcFil] (1.25,0.32) -- (1.25,-0.45) -- (0.45,-0.45);
\pcAppareilRond{(0,-0.45)}{G}
\node[font=\otSemi\fontsize{9}{10}\selectfont] at (-0.6,-0.78) {$-$};
\node[font=\otSemi\fontsize{9}{10}\selectfont] at (0.6,-0.78) {$+$};
\pcNote[west]{(1.61,2.5)}{(2.5,2.6)}{éprouvettes graduées}
\node[pcLegende,anchor=north] at (0,-1.1) {eau rendue conductrice};
\end{tikzpicture}
\par\smallskip\legende{Le montage de l'électrolyse de l'eau ; chaque éprouvette, graduée, recueille le gaz formé à une électrode.}
